Modern multi-domain joint operations present a planning environment in which the battle-space situation evolves rapidly, capabilities must be balanced across competing mission objectives, and decision cycles are measured in seconds rather than minutes. Traditional manual course-of-action development and prompt-only large language model (LLM) assistants fall short on three practical fronts: aligning generated plans with tactical constraints, resisting the variance drift of autoregressive regeneration, and delivering a plan early enough to inform a command-and-control (C2) decision. This paper presents a dual-system enhanced LLM architecture that addresses these constraints as an integrated, evidence-grounded planning loop.

The architecture couples three mechanism families into a single traceable pipeline. A Memento-style case-retrieval module reuses historical operation experience as positive/negative priors; a Hope adaptive controller projects scenario attributes (terrain, weather, electronic-warfare threat, time pressure) onto a bounded capability-weight state through a stochastic, bottleneck-aware update; and a reflection loop converts simulator feedback into local prompt repairs. The pipeline generates structured joint-operation plans, evaluates them in a lightweight closed-loop environment, and exports the selected plan into Department of Defense Architecture Framework (DoDAF) OV-5b, OV-6c, and Command and Control Systems -- Simulation Systems Interoperation (C2SIM) artifacts. On a reference scenario with 10 iterations and 50 Monte Carlo runs, the full loop improves mission success from 62.75 to 65.76 and overall effectiveness from 73.84 to 78.75 over the pure-LLM baseline, and a five-seed replication confirms the gain is statistically significant (paired $t(4)=3.75$, $p=0.02$, $d=1.68$). Across five scene-out generalization scenarios the full pipeline outperforms the pure baseline in four of five, increasing average mission success from 65.01 to 66.48.

For real-time command-and-control use, the architecture adds a decoupled dual-system response path. A rule-based emergency tier produces a structurally complete plan and its DoDAF/C2SIM export in milliseconds (5.6\,ms on the reference hardware, 1.4\,ms warm) with no LLM call, and a bottleneck-driven delta-patching tier refines that anchor to 96.4\% of the full-loop mission-success quality within 22.8\,s, against 365\,s for the complete loop. This yields a three-tier latency profile---millisecond emergency plan, tens-of-seconds agile tactical refinement, minutes-level full optimization---and the delta-patching mechanism preserves the high-quality anchor that full regeneration destroys. The pipeline does not replace dedicated physics-based simulation engines (e.g., OneSAF, FLAMES); instead, it emits IEEE 1516/SISO C2SIM-compatible outputs and acts as a standard-compliant planning hub that bridges LLM-generated courses of action into existing C2 systems and distributed simulation federations. The reported evidence is confined to this internal evaluation sandbox; the paper does not claim external validation against a calibrated field simulator, and the reported confidence intervals reflect within-seed Monte Carlo variance under the fixed seed rather than across-seed variance.
